\documentclass[11pt]{article}
\usepackage{amsmath,amssymb,amsthm}
\usepackage{geometry}
\usepackage{booktabs}
\usepackage{graphicx}
\usepackage{hyperref}
\usepackage{enumerate}
\usepackage{listings}
\usepackage{xcolor}

\geometry{letterpaper, margin=1in}

\lstset{
  basicstyle=\ttfamily\small,
  keywordstyle=\color{blue},
  commentstyle=\color{gray},
  stringstyle=\color{red},
  breaklines=true,
  columns=fullflexible,
  keepspaces=true,
  language=C++
}

\title{\textbf{gas2 Module --- Thermodynamic Potentials and Phase Equilibrium} \\
\Large Theoretical Foundation Document}
\author{VSEPR-SIM Development}
\date{Version 1.0.0 - 4.X Beta --- April 2026}

\begin{document}
\maketitle
\tableofcontents
\newpage

% ============================================================================
\section{Scope and Purpose}
% ============================================================================

This document establishes the theoretical foundation for the thermodynamic
potential infrastructure in the gas2 module.  The goal is to compute
the four fundamental potentials
\[
  A(T,V,N), \quad G(T,P,N), \quad \mu_i(T,P,\mathbf{x}), \quad S(T,V,N)
\]
from a partition-function basis, then derive phase boundaries from rigorous
equilibrium conditions rather than empirical curve fits.

\subsection{Design Principles}

\begin{enumerate}
  \item \textbf{SI internally, Hartree at render time.}
        All thermodynamic quantities are stored in SI (J, Pa, m$^3$, K).
        Hartree conversion is applied at the display layer only:
        \[
          E_\mathrm{Eh} = \frac{E_\mathrm{J}}{E_h}, \qquad
          E_h = 4.3597447222071 \times 10^{-18}\;\mathrm{J}
        \]

  \item \textbf{Anti-black-box.}
        Every intermediate---partition function, fugacity coefficient,
        chemical potential, activity coefficient---is stored in the result
        struct and available for inspection.

  \item \textbf{Deterministic.}
        Given $(T, P, N, \text{species})$, outputs are reproducible
        with no hidden stochastic state.
\end{enumerate}


% ============================================================================
\section{Partition Function and Free Energy}
% ============================================================================

\subsection{Classical Partition Function}

For an $N$-particle system with Hamiltonian $H(\mathbf{q},\mathbf{p})$,
the classical canonical partition function is
\begin{equation}
  Z(T,V,N) = \frac{1}{N!\,h^{3N}}
    \int e^{-\beta H(\mathbf{q},\mathbf{p})}\,d\mathbf{q}\,d\mathbf{p},
    \qquad \beta = \frac{1}{k_B T}.
  \label{eq:partition_function}
\end{equation}

\subsection{Ideal Gas Partition Function}

For a single-component ideal gas of $N$ molecules with molar mass $M$,
the translational partition function factorises:
\begin{equation}
  Z_\mathrm{trans} = \frac{V^N}{N!\,\Lambda^{3N}},
  \qquad
  \Lambda = \sqrt{\frac{h^2}{2\pi m k_B T}}
  \label{eq:thermal_wavelength}
\end{equation}
where $\Lambda$ is the thermal de~Broglie wavelength and $m = M/N_A$.

Internal degrees of freedom contribute multiplicatively:
\begin{equation}
  Z = Z_\mathrm{trans} \cdot Z_\mathrm{rot}^N \cdot Z_\mathrm{vib}^N
      \cdot Z_\mathrm{elec}^N.
\end{equation}

For a rigid rotor:
\begin{equation}
  Z_\mathrm{rot} =
  \begin{cases}
    1, & \text{monoatomic} \\
    \dfrac{T}{\sigma_\mathrm{rot}\,\Theta_\mathrm{rot}}, &
      \text{linear},\;\Theta_\mathrm{rot} = \dfrac{\hbar^2}{2Ik_B} \\[6pt]
    \dfrac{\sqrt{\pi}}{\sigma_\mathrm{rot}}
      \left(\dfrac{T}{\Theta_A}\right)^{1/2}
      \left(\dfrac{T}{\Theta_B}\right)^{1/2}
      \left(\dfrac{T}{\Theta_C}\right)^{1/2}, &
      \text{nonlinear}
  \end{cases}
\end{equation}

For a harmonic oscillator mode $k$:
\begin{equation}
  Z_{\mathrm{vib},k} = \frac{e^{-\Theta_{v,k}/(2T)}}{1 - e^{-\Theta_{v,k}/T}},
  \qquad \Theta_{v,k} = \frac{h\nu_k}{k_B}.
\end{equation}


% ============================================================================
\section{Thermodynamic Potentials}
% ============================================================================

\subsection{Helmholtz Free Energy $A(T,V,N)$}

The Helmholtz free energy follows directly from the partition function:
\begin{equation}
  A(T,V,N) = -k_B T \ln Z(T,V,N).
  \label{eq:helmholtz}
\end{equation}

For the ideal gas:
\begin{equation}
  A^\mathrm{ig}(T,V,N) = -Nk_BT
    \left[\ln\!\left(\frac{V}{N\Lambda^3}\right) + 1\right]
    - Nk_BT \ln Z_\mathrm{int}(T).
  \label{eq:helmholtz_ig}
\end{equation}

For a real gas described by a cubic equation of state (Van der Waals, RK, PR),
the departure function is
\begin{equation}
  A^\mathrm{dep}(T,V,N) = A(T,V,N) - A^\mathrm{ig}(T,V,N)
    = -\int_\infty^V \left(P - \frac{NRT}{V'}\right) dV'.
  \label{eq:helmholtz_departure}
\end{equation}

\paragraph{Van der Waals departure:}
\begin{equation}
  A^\mathrm{dep}_\mathrm{VdW} = -NRT\ln\!\left(\frac{V - Nb}{V}\right)
    - \frac{N^2 a}{V}.
  \label{eq:helmholtz_vdw}
\end{equation}

\paragraph{Redlich--Kwong departure:}
\begin{equation}
  A^\mathrm{dep}_\mathrm{RK} = -NRT\ln\!\left(\frac{V - Nb}{V}\right)
    - \frac{Na}{b\sqrt{T}}\ln\!\left(\frac{V + Nb}{V}\right).
  \label{eq:helmholtz_rk}
\end{equation}

\subsection{Gibbs Free Energy $G(T,P,N)$}

The Gibbs free energy is obtained via Legendre transform:
\begin{equation}
  G(T,P,N) = A(T,V,N) + PV.
  \label{eq:gibbs}
\end{equation}

In practice, we solve for $V(T,P)$ from the EOS, then evaluate:
\begin{equation}
  G = A^\mathrm{ig}(T,V,N) + A^\mathrm{dep}(T,V,N) + PV.
  \label{eq:gibbs_eval}
\end{equation}

\subsection{Entropy}

\begin{equation}
  S(T,V,N) = -\left(\frac{\partial A}{\partial T}\right)_{V,N}.
  \label{eq:entropy}
\end{equation}

For the ideal gas (Sackur--Tetrode):
\begin{equation}
  S^\mathrm{ig} = Nk_B\left[\frac{5}{2}
    + \ln\!\left(\frac{V}{N\Lambda^3}\right)\right]
    + S_\mathrm{int}(T).
  \label{eq:sackur_tetrode}
\end{equation}


% ============================================================================
\section{Chemical Potential and Phase Equilibrium}
% ============================================================================

\subsection{Chemical Potential $\mu_i(T,P,\mathbf{x})$}

For component $i$ in a mixture:
\begin{equation}
  \mu_i(T,P,\mathbf{x}) = \mu_i^\circ(T)
    + RT\ln\!\left(\frac{f_i}{f_i^\circ}\right),
  \label{eq:chemical_potential}
\end{equation}
where $f_i$ is the fugacity of component $i$ and $f_i^\circ$ is the
reference fugacity (typically 1~bar for gases).

The fugacity coefficient from a cubic EOS:
\begin{equation}
  \ln \phi_i = \frac{1}{RT}\int_V^\infty
    \left[\left(\frac{\partial P}{\partial N_i}\right)_{T,V,N_{j\neq i}}
    - \frac{RT}{V}\right] dV - \ln Z.
  \label{eq:fugacity}
\end{equation}

\subsection{Phase Equilibrium Conditions}

At phase equilibrium between phases $\alpha$ and $\beta$:
\begin{align}
  T^\alpha &= T^\beta, \label{eq:thermal_eq} \\
  P^\alpha &= P^\beta, \label{eq:mechanical_eq} \\
  \mu_i^\alpha(T,P,\mathbf{x}^\alpha) &=
    \mu_i^\beta(T,P,\mathbf{x}^\beta)
    \quad \forall\;i.
  \label{eq:chemical_eq}
\end{align}

For a pure substance, this reduces to finding $(T,P)$ where the
Gibbs free energy of the two phases is equal:
\begin{equation}
  G^\alpha(T,P) = G^\beta(T,P).
  \label{eq:gibbs_equality}
\end{equation}

\subsection{Maxwell Construction (Equal-Area Rule)}

For a cubic EOS at subcritical temperatures, the pressure--volume
isotherm exhibits a van der Waals loop.  The saturation pressure
$P_\mathrm{sat}$ is determined by the equal-area condition:
\begin{equation}
  \int_{V_\mathrm{liq}}^{V_\mathrm{vap}}
    \left[P(V) - P_\mathrm{sat}\right] dV = 0,
  \label{eq:maxwell_construction}
\end{equation}
which is equivalent to requiring $G^\mathrm{liq} = G^\mathrm{vap}$.

This is implemented as a root-finding problem: find $P_\mathrm{sat}$
such that the fugacities of the liquid and vapour roots are equal.


% ============================================================================
\section{Intermolecular Potential Decomposition}
% ============================================================================

The total potential energy of a molecular system is decomposed as:
\begin{equation}
  U_\mathrm{total} = U_\mathrm{bond} + U_\mathrm{angle}
    + U_\mathrm{torsion} + U_\mathrm{vdW}
    + U_\mathrm{Coul} + U_\mathrm{pol}
    + U_\mathrm{many\text{-}body}.
  \label{eq:potential_decomposition}
\end{equation}

\subsection{Bonded Interactions}

\paragraph{Bond stretching (harmonic):}
\begin{equation}
  U_\mathrm{bond} = \sum_\mathrm{bonds} \frac{1}{2}\,k_r\,(r - r_0)^2.
\end{equation}

\paragraph{Angle bending (harmonic):}
\begin{equation}
  U_\mathrm{angle} = \sum_\mathrm{angles}
    \frac{1}{2}\,k_\theta\,(\theta - \theta_0)^2.
\end{equation}

\paragraph{Torsional rotation:}
\begin{equation}
  U_\mathrm{torsion} = \sum_\mathrm{dihedrals}
    \frac{V_n}{2}\left[1 + \cos(n\phi - \phi_0)\right].
\end{equation}

\subsection{Non-bonded Interactions}

\paragraph{Van der Waals (Lennard-Jones 12-6):}
\begin{equation}
  U_\mathrm{vdW} = \sum_{i<j} 4\varepsilon_{ij}
    \left[\left(\frac{\sigma_{ij}}{r_{ij}}\right)^{12}
    - \left(\frac{\sigma_{ij}}{r_{ij}}\right)^{6}\right].
\end{equation}

\paragraph{Coulomb electrostatics:}
\begin{equation}
  U_\mathrm{Coul} = \sum_{i<j} \frac{k_e\,q_i\,q_j}{r_{ij}}.
\end{equation}

\paragraph{Polarization (induced dipole, SCF):}
\begin{equation}
  U_\mathrm{pol} = -\frac{1}{2} \sum_i \boldsymbol{\mu}_i^\mathrm{ind}
    \cdot \mathbf{E}_i^\mathrm{perm},
\end{equation}
where $\boldsymbol{\mu}_i^\mathrm{ind} = \alpha_i \mathbf{E}_i^\mathrm{total}$
and the total field is solved self-consistently (see \texttt{polarization\_scf.cpp}).

\paragraph{Many-body corrections:}
\begin{equation}
  U_\mathrm{many\text{-}body} = U_\mathrm{3\text{-}body}
    + U_\mathrm{association} + \cdots
\end{equation}
These include Axilrod--Teller--Muto triple-dipole terms and
hydrogen-bonding association models where applicable.


% ============================================================================
\section{Coarse-Grained Free-Energy Functional}
% ============================================================================

Near a phase transition, the relevant physics is captured by a
density order parameter:
\begin{equation}
  \phi(\mathbf{r}) = \rho(\mathbf{r}) - \rho_c,
  \label{eq:order_parameter}
\end{equation}
where $\rho_c$ is the critical density.

The Landau--Ginzburg free-energy functional:
\begin{equation}
  F[\phi] = \int \left[
    a\,\phi^2 + b\,\phi^4 + \kappa\,(\nabla\phi)^2 + \cdots
  \right] dV,
  \label{eq:landau_ginzburg}
\end{equation}
where:
\begin{itemize}
  \item $a(T) = a_0(T - T_c)$ changes sign at the critical point,
  \item $b > 0$ ensures thermodynamic stability,
  \item $\kappa > 0$ penalises density gradients (interface cost).
\end{itemize}

The equilibrium density profile minimises $F[\phi]$, yielding the
Euler--Lagrange equation:
\begin{equation}
  2a\,\phi + 4b\,\phi^3 - 2\kappa\,\nabla^2\phi = 0.
  \label{eq:euler_lagrange}
\end{equation}


% ============================================================================
\section{Monitoring Infrastructure}
% ============================================================================

The gas2 module exposes a steady-state monitoring window displaying
the 8-component potential decomposition and free-energy functional
in real time.

\subsection{Potential Decomposition Display}

Eight channels of $U$, each updated per evaluation cycle:
\begin{center}
\begin{tabular}{clll}
  \toprule
  Channel & Symbol & Description & Units \\
  \midrule
  1 & $U_\mathrm{bond}$      & Bond stretching        & kJ/mol \\
  2 & $U_\mathrm{angle}$     & Angle bending          & kJ/mol \\
  3 & $U_\mathrm{torsion}$   & Dihedral rotation      & kJ/mol \\
  4 & $U_\mathrm{vdW}$       & Van der Waals          & kJ/mol \\
  5 & $U_\mathrm{Coul}$      & Coulomb electrostatics & kJ/mol \\
  6 & $U_\mathrm{pol}$       & Polarization           & kJ/mol \\
  7 & $U_\mathrm{many}$      & Many-body corrections  & kJ/mol \\
  8 & $U_\mathrm{total}$     & Total potential         & kJ/mol \\
  \bottomrule
\end{tabular}
\end{center}

\subsection{Free-Energy Functional Display}

The free-energy functional $F[\phi]$ is evaluated on a 1D density
profile $\phi(x)$ across the simulation domain, with coefficients
$a$, $b$, $\kappa$ derived from the EOS:
\begin{equation}
  a = \frac{1}{2}\left(\frac{\partial^2 A}{\partial V^2}\right)_{T,N},
  \qquad
  b \sim \frac{1}{4!}\left(\frac{\partial^4 A}{\partial V^4}\right)_{T,N},
  \qquad
  \kappa \propto \sigma^2 \varepsilon,
\end{equation}
where $\sigma$ and $\varepsilon$ are the LJ parameters for the species.


% ============================================================================
\section{Implementation Mapping}
% ============================================================================

\begin{center}
\begin{tabular}{ll}
  \toprule
  \textbf{Theory} & \textbf{Code Location} \\
  \midrule
  $Z(T,V,N)$, $\Lambda$
    & \texttt{include/gas2/gas2\_thermo.hpp} \\
  $A(T,V,N)$, $A^\mathrm{dep}$
    & \texttt{include/gas2/gas2\_thermo.hpp} \\
  $G(T,P,N)$
    & \texttt{include/gas2/gas2\_thermo.hpp} \\
  $\mu_i$, $\phi_i$, $f_i$
    & \texttt{include/gas2/gas2\_thermo.hpp} \\
  $P_\mathrm{sat}$ (Maxwell construction)
    & \texttt{include/gas2/gas2\_thermo.hpp} \\
  $U$ decomposition
    & \texttt{include/gas2/gas2\_potential.hpp} \\
  $F[\phi]$ functional
    & \texttt{include/gas2/gas2\_potential.hpp} \\
  Monitoring window
    & \texttt{tools/gas2\_monitor.py} \\
  Tests
    & \texttt{tests/test\_gas2.cpp} \\
  \bottomrule
\end{tabular}
\end{center}


\end{document}
